\documentclass[english, hyperref={colorlinks,citecolor=blue,linkcolor=blue,urlcolor=blue}, handout]{beamer}

\mode<presentation>
{
  \usetheme{Madrid}
  \usecolortheme{wolverine}
  \usefonttheme{serif}
  \setbeamercovered{transparent}
}

\usepackage[english]{babel}
\usepackage[utf8]{inputenc}
\usepackage{amsmath, amssymb}
\usepackage{graphicx}
\usepackage{booktabs}
\usepackage{xcolor}
\usepackage{url}
\usepackage{colortbl}

% ---------------------------------------------------------------
% Meta
% ---------------------------------------------------------------
\title[KD for On-Device Classification]{%
  Knowledge Distillation for Efficient\\On-Device Image Classification%
}

\author[Saad $\cdot$ Hafeez $\cdot$ Verda]{%
  Saad Ali \texttt{25280011}\\[0.2em]
  Abdul Hafeez \texttt{25280098}\\[0.2em]
  Verda Batool \texttt{24280065}%
}

\institute[LUMS]{Lahore University of Management Sciences}
\date{25 May 2026}

% ===============================================================
\begin{document}
% ===============================================================

% ---------------------------------------------------------------
% Slide 1: Title
% ---------------------------------------------------------------
\begin{frame}
  \titlepage
\end{frame}

% ---------------------------------------------------------------
% Slide 2: Problem Statement
% ---------------------------------------------------------------
\begin{frame}{Problem Statement}
  \begin{columns}[t]
    % ---- left ----
    \begin{column}{0.50\textwidth}
      \textbf{The deployment gap}
      \begin{itemize}\footnotesize
        \item Large CNNs (ResNet-110: \textbf{1.7M params}) achieve high accuracy.
        \item But edge devices have hard limits on memory, latency, and power
        \item Model Compression for NN introduced in 2006
        \item Train a compact \emph{student} to mimic a large pretrained \emph{teacher}
      \end{itemize}
      \vspace{0.35em}
      \textbf{Student trained on Hard Labels}
      \begin{itemize}\small
        \item CCE gradient: $\partial L/\partial z_k^S = \hat{y}_k^S - y_k$
        \item Non-target classes always have $y_k{=}0$, pushing all wrong logits to $-\infty$ regardless of class similarity
      \end{itemize}
    \end{column}
    % ---- right ----
    \begin{column}{0.46\textwidth}
      \textbf{Teacher's soft output encodes \emph{dark knowledge}}
      \vspace{0.3em}

      \begin{center}\tiny
        \begin{tabular}{lcc}
          \toprule
          Class & Hard label & Soft ($\tau{=}4$)\\
          \midrule
          \texttt{dog}   & \textbf{1.00} & \textbf{0.71}\\
          \texttt{wolf}  & 0.00 & \textit{0.14}\\
          \texttt{fox}   & 0.00 & \textit{0.09}\\
          \texttt{cat}   & 0.00 & \textit{0.04}\\
          other 96       & 0.00 & \textit{0.02}\\
          \bottomrule
        \end{tabular}
      \end{center}
      \vspace{0.1em}
      \begin{block}{\small Dark Knowledge \textnormal{\small[Hinton 2015]}}\small
        Nonzero probabilities on wrong classes encode inter-class similarity that hard labels discard.
      \end{block}
    \end{column}
  \end{columns}
\end{frame}

% ---------------------------------------------------------------
% Slide 4: Methodology
% ---------------------------------------------------------------
\begin{frame}{Methodology: Response-Based Knowledge Distillation}
  \begin{columns}[t]
    % ---- left ----
    \begin{column}{0.50\textwidth}
      \textbf{Temperature-scaled softmax}
      \[
        \hat{y}_k^{T,\tau} = \frac{\exp(z_k^T/\tau)}{\textstyle\sum_j \exp(z_j^T/\tau)}
      \]
      {\small $\tau{>}1$ flattens distribution $\Rightarrow$ amplifies non-target gradients}

      \vspace{0.4em}
      \textbf{Distillation objective}
      \[
        \boxed{L_{\mathrm{KD}}
          = \alpha\,L_{\mathrm{hard}}
          + (1{-}\alpha)\,\tau^2 L_{\mathrm{soft}}}
      \]
      {\small $\tau^2$ compensates gradient shrinkage: $\partial L_{\mathrm{soft}}/\partial z_k^S \propto 1/\tau^2$}

      \vspace{0.4em}
      \textbf{Hyperparameter sweeps}
      \begin{itemize}\small
        \item Temperature: $\tau \in \{2, 4, 8, 16\}$
        \item Loss weight: $\alpha \in \{0.1, 0.5, 0.9\}$
      \end{itemize}
    \end{column}
    % ---- right ----
    \begin{column}{0.47\textwidth}
      \textbf{Capacity study}\\[0.2em]
      \begin{itemize}\small
        \item \textit{Depth}: ResNet-$\{20,32,44,56\}$, width 16
        \item \textit{Width}: ResNet-20, multiplier $w\in\{0.5,1,2\}$
        \item Shared pivot: ResNet-20 ($w{=}1$)
      \end{itemize}

      \vspace{0.5em}
      \textbf{Training config}
      \begin{center}\tiny
        \begin{tabular}{ll}
          \toprule
          Setting & Value\\
          \midrule
          Optimiser      & SGD + Momentum (0.9)\\
          LR / decay     & 0.1 / ${\times}0.1$ at ep.\ 60, 80\\
          Weight decay   & $5{\times}10^{-4}$\\
          Batch / Epochs & 128 / 100 (student)\\
          \bottomrule
        \end{tabular}
      \end{center}
    \end{column}
  \end{columns}
\end{frame}

% ---------------------------------------------------------------
% Slide 3: Dataset & Experimental Setup
% ---------------------------------------------------------------
\begin{frame}{Dataset \& Experimental Setup}
  \begin{columns}[t]
    % ---- left ----
    \begin{column}{0.42\textwidth}
      \textbf{CIFAR-100}
      \begin{itemize}\small
        \item 60{,}000 images, $32{\times}32$, 100 classes
        \item 50k train / 10k test (500 per class)
      \end{itemize}
      \vspace{0.4em}
      \textbf{Models (best-checkpoint)}
      \begin{center}\tiny
        \begin{tabular}{lcc}
          \toprule
          & Params & Top-1\\
          \midrule
          Teacher (RN-110) & 1.74M & 73.48\%\\
          \rowcolor{yellow!30}
          Baseline (RN-20) & 0.28M & 67.63\%\\
          \bottomrule
        \end{tabular}
      \end{center}
      \vspace{0.2em}
      {\small\textbf{Gap = 5.85 pp} = distillation headroom}
    \end{column}
    % ---- right ----
    \begin{column}{0.55\textwidth}
      \textbf{Baseline ResNet-20 training on CIFAR-100 (100 epochs).}
      \vspace{0.1em}

      \includegraphics[width=\linewidth, height=4.2cm, keepaspectratio]%
        {plots/baseline/resnet20/training_curves.png}

      \vspace{0.1em}
      {\footnotesize LR decays at epochs 60 \& 80 produce the visible accuracy jumps. Train$-$test gap $\approx\!12.5$ pp indicates mild overfitting.}
    \end{column}
  \end{columns}
\end{frame}


% ---------------------------------------------------------------
% Slide 5: Results
% ---------------------------------------------------------------
\begin{frame}{Results: Hyperparameter Sweeps \& Capacity Scaling}
  \begin{columns}[t]
    % ---- left: tables ----
    \begin{column}{0.38\textwidth}
      \textbf{Temperature} ($\alpha{=}0.5$)
      \begin{center}\scriptsize
        \begin{tabular}{ccc}
          \toprule
          $\tau$ & Top-1 & $\Delta$\\
          \midrule
          2  & 69.49\% & $+1.86$\\
          4  & 69.64\% & $+2.01$\\
          \rowcolor{green!25}
          \textbf{8}  & \textbf{69.98\%} & $\mathbf{+2.35}$\\
          16 & 69.73\% & $+2.10$\\
          \bottomrule
        \end{tabular}
      \end{center}
      {\scriptsize Inverted-U; $\tau^*{=}8$ (Hinton: $\tau{\approx}4$)}

      \vspace{0.5em}
      \textbf{Loss weight} ($\tau{=}8$)
      \begin{center}\scriptsize
        \begin{tabular}{ccc}
          \toprule
          $\alpha$ & Top-1 & $\Delta$\\
          \midrule
          0.1 & 69.74\% & $+2.11$\\
          \rowcolor{green!25}
          \textbf{0.5} & \textbf{69.98\%} & $\mathbf{+2.35}$\\
          0.9 & 69.07\% & $+1.44$\\
          \bottomrule
        \end{tabular}
      \end{center}
      {\scriptsize $\alpha{=}0.9$ recovers only 62\% of peak gain; soft labels provide genuinely new information}
    \end{column}
    % ---- right: capacity figure ----
    \begin{column}{0.60\textwidth}
      \textbf{Capacity scaling} ($\tau{=}8$, $\alpha{=}0.5$)

      \includegraphics[width=\linewidth, height=4.5cm, keepaspectratio]%
        {plots/analysis/capacity.png}

      {\footnotesize \textbf{3 of 5 students exceed teacher (73.48\%).} Best: ResNet-20 $w{=}2$ reaches \textbf{75.43\%} (+7.80 pp), still 1.6$\times$ smaller. Low-capacity failure: $w{=}0.5$ falls \textbf{10.37 pp below baseline}.}
    \end{column}
  \end{columns}
\end{frame}

% ---------------------------------------------------------------
% Slide 6: Dark Knowledge & Calibration
% ---------------------------------------------------------------
\begin{frame}{Dark Knowledge \& Calibration Analysis}
  \begin{columns}[t]
    % ---- left: dark knowledge ----
    \begin{column}{0.52\textwidth}
      \textbf{Why soft labels help: dark knowledge}

      \includegraphics[width=\linewidth, height=3.3cm, keepaspectratio]%
        {plots/analysis/dark_knowledge.png}

      {\footnotesize Each row: image $|$ teacher $\tau{=}1$ $|$ teacher $\tau{=}4$ $|$ hard-label student $|$ distilled student.
      At $\tau{=}4$, probability spreads to similar classes; distilled student mimics this structure while hard-label student does not.}
    \end{column}
    % ---- right: calibration ----
    \begin{column}{0.45\textwidth}
      \textbf{Side-effect: calibration degrades}
      \begin{center}\tiny
        \begin{tabular}{lcc}
          \toprule
          Model & Top-1 & ECE$\downarrow$\\
          \midrule
          Teacher           & 73.48\% & 12.54\%\\
          \rowcolor{green!20}
          Hard-label (base) & 67.63\% & \textbf{6.23\%}\\
          Distilled ($\tau{=}8$) & 69.98\% & 10.32\%\\
          \rowcolor{yellow!30}
          Distilled ($w{=}2$) & \textbf{75.43\%} & 9.97\%\\
          \bottomrule
        \end{tabular}
      \end{center}

      \includegraphics[width=\linewidth, height=2.4cm, keepaspectratio]%
        {plots/analysis/calibration.png}

      {\footnotesize KL term transfers teacher's \textbf{overconfidence profile}, not just the argmax. Hard-label student is best calibrated despite lowest Top-1. Post-hoc calibration required for probability-sensitive applications.}
    \end{column}
  \end{columns}
\end{frame}

% ===============================================================
\end{document}
